\documentclass{article}
\usepackage{amsmath, amssymb, amsthm}
\usepackage{hyperref}
\usepackage{geometry}
\geometry{margin=1in}

\title{Erd\H{o}s Problem \#675}
\date{Last updated: 2025-08-31}
\author{Source: erdosproblems.com}

\begin{document}
\maketitle

\noindent\textbf{Status:} OPEN \quad \textbf{No prize}

\noindent\textbf{Tags:} number theory

\medskip

\noindent\textbf{Problem Statement:}

We say that $A\subset \mathbb{N}$ has the translation property if, for every $n$, there exists some integer $t_n\geq 1$ such that, for all $1\leq a\leq n$,\[a\in A\quad\textrm{ if and only if }\quad a+t_n\in A.\]{UL} {LI}Does the set of the sums of two squares have the translation property?{/LI} {LI}If we partition all primes into $P\sqcup Q$, such that each set contains $\gg x/\log x$ many primes $\leq x$ for all large $x$, then can the set of integers only divisible by primes from $P$ have the translation property?{/LI} {LI}If $A$ is the set of squarefree numbers then how fast does the minimal such $t_n$ grow? Is it true that $t_n&#62;\exp(n^c)$ for some constant $c&#62;0$?{/LI} {/UL}




Elementary sieve theory implies that the set of squarefree numbers has the translation property. More generally, Brun's sieve can be used to prove that if $B\subseteq \mathbb{N}$ is a set of pairwise coprime integers with $\sum_{b&#60;x}\frac{1}{b}=o(\log\log x)$ then $A=\{ n: b\nmid n\textrm{ for all }b\in A\}$ has the translation property. Erd\H{o}s did not know what happens if the condition on $\sum_{b&#60;x}\frac{1}{b}$ is weakened or dropped altogether.


\bigskip
\noindent\small{Source: \url{https://www.erdosproblems.com/675}}
\end{document}
